\section{Nothing Here Is Authorized}
\label{sec:ch15-nothing-is-authorized}
\index{authorization!the state before any of it|(}%

Check out tag \code{ch14}, start the graph, and read somebody's email address:

\begin{minted}[fontsize=\footnotesize,breaklines]{graphql}
{ products { reviews(first: 1) { nodes { author { displayName email } } } } }
\end{minted}

Twenty-five products, an author on nearly every one, and an address for each.
There is no token in that request and no place to put one. Ask Ordering for
somebody's purchase history and it is the same story, one argument and no
question about who is asking.

I want to be precise about what is wrong here, because \enquote{the graph has no
authorization} is a description of a to-do item rather than of a defect. The
defect is narrower and it is worth naming: \code{Customer.email} is a field
Accounts publishes for the convenience of a graph that reaches customers through
their reviews, and a review is public. Nobody decided that an address should
follow a display name into a product page. It followed because the two are
properties of one type, and a type is the unit federation composes.

That is the shape of the problem for the rest of the chapter. The units a
federated graph is assembled from are types and fields; the units authorization
cares about are rows and people.

\subsection*{Three places to put a rule}

\index{authorization!where a rule can live}%
Before any code, the decision that everything else follows from. A request to
this graph passes through at least two processes, and it can be refused in three
places, which are not interchangeable.

\index{authorization!at the router}%
The router can refuse a field from the token alone. It has the claims, it has
the operation, and the composed config tells it which fields need what. It does
this before planning a fetch, so a refusal costs nothing. What it cannot do is
know whose data it was about to fetch, because it has not fetched it yet.

\index{authorization!at the subgraph}%
The subgraph can refuse the same field inside its own process, using ASP.NET
Core's authorization and HotChocolate's \code{[Authorize]}. That looks redundant
next to the first one and is not, for a reason that is a property of this
deployment rather than of federation: all seven Mosaic services listen on a host
port. A client that skips the router is one \code{curl} away, and every rule the
router was holding evaporates.

\index{authorization!in the resolver}%
The resolver can refuse this row to this caller. An order belongs to one
customer, and the question is not whether you hold a scope but whether you are
that customer. Answering it needs the row, so it happens after the fetch, in the
service that owns the data, and no directive anywhere can advertise that it
happened.

Figure~\ref{fig:ch15-three-layers} is the three of them on one path, including
the arrow that makes the second one necessary.

\begin{figure}[htbp]
  \centering
  % Three places a rule can live, on one path. Referenced from chapter 15. Bare
% tikzpicture; the figure environment, caption and label live at the call site.
%
% The point of the picture is that the three checkpoints are not
% interchangeable: each one can say something the others cannot, and each one
% sits where it does because that is where the thing it asks about is. The
% dashed arrow is what makes the middle one necessary rather than redundant.
\begin{tikzpicture}[
    font=\small,
    box/.style={draw, rounded corners=2pt, minimum width=18mm,
                minimum height=7mm, align=center, font=\scriptsize},
    gate/.style={draw, rounded corners=2pt, minimum width=24mm,
                 minimum height=8mm, align=center, font=\scriptsize,
                 fill=black!8},
    flow/.style={-{Stealth[length=2mm]}},
    bypass/.style={-{Stealth[length=2mm]}, dashed},
    hangs/.style={gray},
    note/.style={font=\scriptsize, gray, align=center, inner sep=1pt}
  ]

  % -- the path a request takes ---------------------------------------------
  \node[box] (client)   at (0.7,0)  {client};
  \node[box] (router)   at (4.0,0)  {router};
  \node[box] (ordering) at (7.6,0)  {\code{ordering}};
  \node[box] (row)      at (11.1,0) {the row};

  \draw[flow] (client)   -- (router);
  \draw[flow] (router)   -- (ordering);
  \draw[flow] (ordering) -- (row);

  % -- the rule each one is able to apply ------------------------------------
  \node[gate] (r1) at (4.0,2.0)  {\code{@authenticated}\\\code{@requiresScopes}};
  \node[gate] (r2) at (7.6,2.0)  {\code{[Authorize]}};
  \node[gate] (r3) at (11.1,2.0) {\code{CustomerId}\\\code{== the subject}};

  \draw[hangs] (r1) -- (router);
  \draw[hangs] (r2) -- (ordering);
  \draw[hangs] (r3) -- (row);

  \node[note, anchor=south] at (4.0,2.6)  {read from the token,\\before any fetch};
  \node[note, anchor=south] at (7.6,2.6)  {this process only,\\and no router sees it};
  \node[note, anchor=south] at (11.1,2.6) {needs the row,\\so no schema can say it};

  % -- and the door the first rule does not cover ----------------------------
  \draw[bypass] (client.south) to[out=-45, in=-135] (ordering.south);
  \node[note, anchor=north] at (4.15,-2.1)
    {every service listens on a host port, so this request happens\\
     and only the middle rule is left to refuse it};

\end{tikzpicture}

  \caption{One field, three rules, three places. The router answers from the
  token before it fetches; the subgraph answers in its own process, which is the
  only rule left for a caller who never went through the router; the resolver
  answers with the row in hand, which is the only place the question \enquote{is
  this yours} can be asked at all. Each layer can express something the ones
  above it cannot, and nothing in the build checks that the three agree.}
  \label{fig:ch15-three-layers}
\end{figure}

\subsection*{What Mosaic decides}

\index{Mosaic!the authorization model}%
Four fields get guarded, and they were chosen to cover the three layers rather
than to be a complete policy. \code{Customer.email} needs a caller who is
somebody. \code{Query.customerById} needs a scope, because looking a customer up
by identifier is a different act from following a review to its author.
\code{Query.ordersByCustomer} needs a scope and needs you to be the customer in
the argument. \code{submitReview} needs a caller and needs the review to be in
your own name, and its scope is checked only inside Reviews rather than at the
router, which is a difference section~\ref{sec:ch15-two-vocabularies} is about
and not an oversight.

\index{tokens!what Mosaic puts in one}%
The token is HS256 with a shared secret, minted by
\code{scripts/mint-token.mjs}, and I want to be blunt about what that costs.
A symmetric key means every party that can check a token can also mint one, and
the key is in a public repository. A deployment points the router at a JWKS
endpoint instead and gives the services nothing; the router's configuration
takes either and \code{router/config.yaml} carries the two-line swap as a
comment. Nothing else in the repository changes between them, which is the only
reason I am comfortable spending a symmetric key on a chapter about federation
rather than a chapter about OAuth.

\index{tokens!the subject is a global object identifier}%
One decision inside the token is worth more than it looks. The \code{sub} claim
is a customer's global object identifier, the same base64 string
chapter~\ref{ch:schema-design-that-survives-change} built,
chapter~\ref{ch:the-first-cut-extracting-catalog} first taught a subgraph to
decode, and chapter~\ref{ch:strangling-the-monolith} copied into every service
that needed one. So a service that wants the customer behind a token decodes the
subject with its own copy of \code{CustomerKey}, and there are two more copies at
this tag for exactly that. Chapter~\ref{ch:strangling-the-monolith} argued that the
format of an identifier is a contract and a shared class would turn it into a
build dependency. This adds a party to the contract who is not a Mosaic service
at all: whatever mints tokens has to encode a customer the way Accounts does,
and no compiler will ever check it.

\medskip
So: a token, and somewhere to check it. The router first, because it is the one
that can refuse a field before anything is fetched.

\index{authorization!the state before any of it|)}%
